%%%%%%%%%%%%%%%%%%%%%%%%%%%%%%%%%%%%%%%%%%%%%%%%%%%%%%%%%%%%%%%%%%%%%%
\section{Questions et points de discussion}
%%%%%%%%%%%%%%%%%%%%%%%%%%%%%%%%%%%%%%%%%%%%%%%%%%%%%%%%%%%%%%%%%%%%%%

L'objectif de cette note n'est pas de présenter une méthodologie figée.
Au contraire, plusieurs choix techniques et scientifiques restent ouverts
et mériteraient d'être discutés avec un enseignant-chercheur avant le
démarrage effectif du projet.

Les questions suivantes constituent les principaux points sur lesquels
je souhaiterais bénéficier de votre expérience.

%%%%%%%%%%%%%%%%%%%%%%%%%%%%%%%%%%%%%%%%%%%%%%%%%%%%%%%%%%%%%%%%%%%%%%

\subsection{1. Constitution du corpus}

\begin{discussion}

\begin{itemize}

\item Quel volume de données serait raisonnable pour une première
expérimentation scientifique ?

\item Est-il préférable de collecter un grand nombre de locuteurs avec
peu d'enregistrements chacun, ou un nombre plus réduit de locuteurs avec
davantage d'enregistrements ?

\item Faut-il commencer par des mots isolés, des phrases simples ou
directement des conversations spontanées ?

\item Existe-t-il des méthodologies de collecte que vous recommanderiez
pour une langue peu dotée comme le soussou ?

\end{itemize}

\end{discussion}

%%%%%%%%%%%%%%%%%%%%%%%%%%%%%%%%%%%%%%%%%%%%%%%%%%%%%%%%%%%%%%%%%%%%%%

\subsection{2. Métadonnées}

\begin{discussion}

\begin{itemize}

\item Quelles métadonnées vous semblent indispensables pour garantir la
qualité scientifique du corpus ?

\item Certaines informations devraient-elles être évitées afin de limiter
les biais ou les problèmes liés à la confidentialité ?

\item Existe-t-il un format standard que vous recommanderiez pour
l'organisation des données ?

\end{itemize}

\end{discussion}

%%%%%%%%%%%%%%%%%%%%%%%%%%%%%%%%%%%%%%%%%%%%%%%%%%%%%%%%%%%%%%%%%%%%%%

\subsection{3. Prétraitement des données}

\begin{discussion}

\begin{itemize}

\item Quels formats audio recommandez-vous pour ce type de projet
(16 kHz, 44,1 kHz, PCM, FLAC, WAV...) ?

\item Quels contrôles qualité devraient être réalisés avant toute phase
d'apprentissage automatique ?

\item Quels outils ou bibliothèques vous semblent les plus adaptés pour
cette étape ?

\end{itemize}

\end{discussion}

%%%%%%%%%%%%%%%%%%%%%%%%%%%%%%%%%%%%%%%%%%%%%%%%%%%%%%%%%%%%%%%%%%%%%%

\subsection{4. Modélisation}

\begin{discussion}

\begin{itemize}

\item Quels modèles de représentation de la parole vous sembleraient les
plus adaptés pour un corpus de taille limitée ?

\item Recommanderiez-vous une approche auto-supervisée, supervisée ou
hybride dans les premières expérimentations ?

\item Existe-t-il des travaux ou des modèles récents que vous me
conseilleriez d'étudier avant de commencer cette phase ?

\end{itemize}

\end{discussion}

%%%%%%%%%%%%%%%%%%%%%%%%%%%%%%%%%%%%%%%%%%%%%%%%%%%%%%%%%%%%%%%%%%%%%%

\subsection{5. Vision scientifique}

\begin{discussion}

\begin{itemize}

\item La progression proposée dans ALORA vous paraît-elle cohérente du
point de vue scientifique ?

\item Voyez-vous des risques méthodologiques importants dès cette
première phase ?

\item Certaines étapes devraient-elles être réorganisées ou approfondies
avant le début du développement ?

\item Pensez-vous que cette problématique puisse conduire, à terme, à un
véritable sujet de recherche en traitement automatique de la parole ou
en linguistique computationnelle ?

\end{itemize}

\end{discussion}

%%%%%%%%%%%%%%%%%%%%%%%%%%%%%%%%%%%%%%%%%%%%%%%%%%%%%%%%%%%%%%%%%%%%%%

\subsection{Conclusion}

Je suis pleinement conscient que les hypothèses présentées dans ce
document pourront évoluer au cours du projet.

Mon objectif n'est pas de défendre une solution déjà arrêtée, mais de
construire progressivement une méthodologie rigoureuse en m'appuyant sur
les recommandations d'un enseignant-chercheur et sur les résultats
obtenus au fil des expérimentations.

Je serais très heureux d'avoir l'occasion d'échanger avec vous afin de
recueillir vos remarques, vos critiques ainsi que vos suggestions avant
de commencer cette première phase de développement.